\section{导读：2025 年的主角从模型公司转向应用和 Agent}

本节先定位这期。张小珺在开场说，2024 年 AI 叙事还是大模型公司，到了 2025 年，新的主角变成应用公司和 Agent。明超平是 1995 年出生的 AI 应用创业者，先后在 OnePlus、ByteDance、Moonshot 做产品，第一次做 CEO，发布的第一个产品叫 YouWare。这期的价值，不只是听一个年轻 founder 的个人故事，而是看一个产品型创始人如何把手机、短视频、Kimi、coding、Agent 和 OS 机会串成一套判断。

\begin{importantbox}{本期核心命题}
YouWare 的机会不只是“做一个 Agent 产品”，而是重新定义 AI 时代的创作容器和执行环境。明超平的判断是：基础模型会持续变好，但应用公司必须通过环境、经验、交互和网络效应，把更聪明的模型适配到具体生产需求里。
\end{importantbox}

\begin{knowledgebox}{视觉策略说明}
本视频是固定访谈画面，没有 slides、白板或产品演示。正文只使用封面作为来源识别；正文图像全部为自制概念图，用来解释产品经理三站、体验与数据、bitter lesson、token 消耗、容器/环境、Vibe Coder、AgentRank、OS Agent 和 Ego 对抗。
\end{knowledgebox}

\begin{figure}[H]
\centering
\includegraphics[width=0.96\textwidth]{product-manager-three-stations.png}
\caption{产品经理三站：OnePlus、ByteDance、Moonshot 形成产品判断底盘。自制概念图，依据 00:03:16--01:05:05 对谈内容整理。}
\end{figure}

\begin{knowledgebox}{读图：三段经历分别训练三种产品能力}
OnePlus 训练的是体验和审美，ByteDance 训练的是数据、增长和高频迭代，Moonshot 训练的是模型、AI 产品和前沿不确定性。YouWare 不是凭空出现，而是这三种能力在 Agent 应用窗口里的重新组合。
\end{knowledgebox}

\subsection{为什么这不是一篇普通创业访谈}

本节先把注意力从个人履历拉回产品方法。明超平的故事里有两条线：一条是人如何形成产品判断，另一条是 AI 应用公司如何在模型巨头旁边找到位置。前者包括辩论训练、智能车竞赛、OnePlus 的体验训练和字节的数据训练；后者包括 bitter lesson、token 消耗、容器环境、AgentRank 和 OS Agent。

\begin{knowledgebox}{课堂提示：把人物故事读成方法论}
这期的有效读法不是“年轻 founder 很有个性”，而是“哪些经历让他形成了哪些产品判断”。辩论对应第三方视角，智能车竞赛对应 reward 和动手反馈，OnePlus 对应体验，字节对应数据，Moonshot 对应模型时代的边界感。
\end{knowledgebox}

\subsection{阅读路线}

本节给出读法。访谈前半段讲辩论、OnePlus、字节和 Kimi，容易被误读为个人履历；后半段讲 token、容器、Vibe Coder、AgentRank 和 OS Agent，才是产品判断的核心。正确读法是：前半段解释他为什么这样做产品，后半段解释 AI 应用公司为什么还有机会。

\noindent\begin{tabularx}{\textwidth}{p{0.18\textwidth}p{0.34\textwidth}X}
\toprule
阅读问题 & 访谈中的材料 & 要形成的判断 \\
\midrule
产品直觉从哪来 & 辩论、OnePlus、字节、Kimi & 第三方视角、体验敏感、数据基建和模型窗口共同塑造。 \\
AI 应用怎么判断 & bitter lesson、token 消耗、per token valuation & 不要雕花，围绕模型能力和智能消耗效率重构产品。 \\
Agent 生态怎么来 & AgentRank、网络效应、身份证、信任 & Agent 不是孤立工具，会形成调用、排序和信任网络。 \\
CEO 机制怎么建 & 情绪价值、如履薄冰、对抗 Ego、棋手心态 & 应用创业需要创始人持续校准自己和组织。 \\
\bottomrule
\end{tabularx}

\subsection{本章小结}

EP101 的主线是：AI 应用创业从模型崇拜走向环境构建。YouWare 的案例说明，应用公司不能只等待模型变强，而要决定智能被放在哪个容器里、怎样被消耗、如何形成网络效应。

